%% This is file `medima-template.tex',
%% 
%% Copyright 2018 Elsevier Ltd
%% 
%% This file is part of the 'Elsarticle Bundle'.
%% ---------------------------------------------
%% 
%% It may be distributed under the conditions of the LaTeX Project Public
%% License, either version 1.2 of this license or (at your option) any
%% later version.  The latest version of this license is in
%%    http://www.latex-project.org/lppl.txt
%% and version 1.2 or later is part of all distributions of LaTeX
%% version 1999/12/01 or later.
%% 
%% The list of all files belonging to the 'Elsarticle Bundle' is
%% given in the file `manifest.txt'.
%% 
%% Template article for Elsevier's document class `elsarticle'
%% with harvard style bibliographic references
%%
%% $Id: medima-template.tex 153 2018-12-01 11:38:32Z rishi $
%% $URL: http://lenova.river-valley.com/svn/elsarticle/trunk/medima-template.tex $
%%
%% Use the option review to obtain double line spacing
%\documentclass[times,review,preprint,authoryear]{elsarticle}

%% Use the options `twocolumn,final' to obtain the final layout
%% Use longtitle option to break abstract to multiple pages if overfull.
%% For Review pdf (With double line spacing)
%\documentclass[times,twocolumn,review]{elsarticle}
%% For abstracts longer than one page.
%\documentclass[times,twocolumn,review,longtitle]{elsarticle}
%% For Review pdf without preprint line
%\documentclass[times,twocolumn,review,nopreprintline]{elsarticle}
%% Final pdf
\documentclass[times,onecolumn,final]{elsarticle}
%%
%\documentclass[times,twocolumn,final,longtitle]{elsarticle}
%%


%% Stylefile to load MEDIMA template
\usepackage{medima_}
\usepackage{framed,multirow}

%% The amssymb package provides various useful mathematical symbols
\usepackage{amssymb}
\usepackage{latexsym}

% Following three lines are needed for this document.
% If you are not loading colors or url, then these are
% not required.
\usepackage{url}
\usepackage{xcolor}

\usepackage{hyperref}

\usepackage{multirow}
\usepackage{mathtools}
\usepackage{bbm}
\usepackage{bm}
\usepackage{physics}
\usepackage[normalem]{ulem}
\useunder{\uline}{\ul}{}
\usepackage{adjustbox}

\usepackage{float}

% 

\definecolor{newcolor}{rgb}{.8,.349,.1}

\journal{Medical Image Analysis}

\begin{document}

\verso{Yi Gu \textit{et~al.}}

\begin{frontmatter}

\title{SUPPLEMENTAL MATERIALS \newline Bone mineral density estimation from a plain x-ray image by learning decomposition into projections of bone-segmented computed tomography}

% \author[1,4]{Yi \snm{Gu}\corref{cor1}}
% \ead{gu.yi.gu4@is.naist.jp}
% \cortext[cor1]{Corresponding authors at ICB lab., Division of Information Science, Graduate School of Science and Technology, Nara Institute of Science and Technology, Japan.}
% \author[1]{Yoshito \snm{Otake}\corref{cor1}}
% \ead{otake@is.naist.jp}

% \author[2]{Keisuke \snm{Uemura}}

% \author[1]{Mazen \snm{Soufi}}

% \author[3]{Masaki \snm{Takao}}

% \author[4]{Hugues \snm{Talbot}}

% \author[2]{Seiji \snm{Okada}}

% \author[5]{Nobuhiko \snm{Sugano}}

% \author[1]{Yoshinobu \snm{Sato}\corref{cor1}}
% \ead{yoshi@is.naist.jp}


\end{frontmatter}

%\linenumbers

%% main text

\newcommand{\rmtxt}[1]{\textcolor{red}{\sout{#1}}}
\newcommand{\adtxt}[1]{{#1}}
% \pagebreak
\begin{figure*}[h]
\centering
\includegraphics[width=\textwidth]{315_pose_to_pose_bmd.pdf}
\caption{Correlations of predicted BMD between poses, including std. (standing), sup. (supine), abd. (abduction), and add. (adduction), obtained by the proposed method on the experiments using (blue scatters) 4 poses and (green scatters) 2 poses in training.
Statistical significance was denoted by $p$-value, which was computed using Tukey honest significant differences test, where the difference between two distributions with $p>0.05$ was considered non-significant.}
\label{fig:315_pose_to_pose_bmd}
\end{figure*}

\begin{figure*}[!t]
\centering
\includegraphics[width=\textwidth]{compress_representative.pdf}
\caption{Sampled good, medium, and worse cases with respect to the BMD estimation accuracy in compression experiment, showing the ROI of original x-ray images, compressed x-ray images, ground-truth (GT) DRR, predicted (Pred.) DRR in stage 2, the absolute error (AE) map between pred. and GT DRR.}
\label{fig:compress_representative}
\end{figure*}

\begin{figure*}[!t]
\centering
\includegraphics[width=\textwidth]{315_quant_vs_uncalib.pdf}
\caption{Correlations of predicted BMD by the proposed method between training with QCT and uncalibrated CT grouped by the poses used in training.
Statistical significance was indicated by $p$-values, where $p$ larger than 0.05 were considered non-significant.}
\label{fig:315_quant_vs_uncalib}
\end{figure*}

\begin{figure*}[!t]
\centering
\includegraphics[width=\textwidth]{angle_vs_error_scatters.pdf}
\caption{\adtxt{Scatters of hip angles against BMD prediction error for error source investigation. The Pearson correlation coefficients (PCCs) for all poses, standing pose, supine pose, abduction pose, and adduction pose are shown at the top of each scatter. Though a weak trend (PCC=-0.273) was observed between flexion angle and DXA-BMD prediction error, no strong correlation appeared overall, likely because of the noise on the femoral head deformation caused by diseases in our dataset.}}
\label{fig:angle_vs_error}
\end{figure*}

\begin{figure*}[!t]
\centering
\includegraphics[width=\textwidth]{315_methods_bland_altman.pdf}
\caption{\adtxt{Bland-altman plots for the conventional reg. and proposed decomp. methods on the (left) DXA-BMD and (right) QCT-BMD estimation tasks.}}
\label{fig:methods_bland_altman}
\end{figure*}

\begin{figure*}[!t]
\centering
\includegraphics[width=\textwidth]{315_screenshot.png}
\caption{\adtxt{Screenshot of visualization video of 1206 samples (305 patients), visualizing the hemi-hip X-ray image, the ROI X-ray image (used to train conventional methods), the ROI-PF-DRR, the ROI-PF label., and BMD estimation results.}}
\label{fig:315_screenshot}
\end{figure*}

\begin{figure*}[!t]
\centering
\includegraphics[width=\textwidth]{200_psnr_dc_box.png}
\caption{Quantitative evaluation results of image decomposition for different backbones w/o and w/ hierarchical learning (HL). 
The performance improvements by HL were observed for all the backbones, where the ResNet backbone benefited the most from HL. 
The lowest accuracy among the methods with HL (DAFormer w/ HL) was higher than the best accuracy among the methods without HL (DAFormer w/o HL). 
The comparison between HRNet and HRFormer suggests that HRNet slightly outperformed the HRFormer under no HL, while applying HL boosted both of them and made HRFormer a better choice.}
\label{fig:315_quant_vs_uncalib}
\end{figure*}

\begin{figure*}[!t]
\centering
\includegraphics[width=\textwidth]{200_psnr_dc_test.png}
\caption{Results of the statistical test on PSNR and DC between methods. 
The $p$ values between the methods without HL and with HL were smaller than 0.001 for all backbones, indicating significant differences. 
The test results showed the DAformer, HRNet, and HRFormer had no significant difference when HL was not used, whereas the differences were more pronounced when HL was used.}
\label{fig:315_quant_vs_uncalib}
\end{figure*}

\begin{figure*}[!t]
\centering
\includegraphics[width=\textwidth]{200_video_screenshot.png}
\caption{Screenshot of visualization video of 200-cases, visualizing decomposition results by HRFormer w/o and w/ HL in the descending order of PSNR by HRFormer w/ HL, where the more consistent shape and more detailed inner structure by applying HL were observed visually. 
The absolute error maps suggest the significant error is more likely produced by the upper and bottom cutting edges, which are clinically defined while applying HL made the model better in handling the cutting edges.}
\label{fig:315_quant_vs_uncalib}
\end{figure*}


\end{document}
%%
